We introduced the Saliency Alignment Loss, a parameter-free, architecture-agnostic auxiliary objective that supervises a VLM's token-conditioned attention with region-level grounding annotations. On LLaVA-1.5-7B, only $200$ steps on grounded COCONut-PanCap captions transform attention localization---mean Attention Mass Ratio rises $4.6{\times}$, Average Precision doubles, and Normalized Scanpath Saliency rises sevenfold---while leaving language modeling intact, and the qualitative maps confirm it: mass the base model scattered onto image borders and background is redrawn onto the objects a token names.

Our analysis gives an unusually complete account of \emph{where} this happens. The language model, not the multimodal projector, absorbs the signal; it does so identically whether or not the projector is also trainable, reaching a reproducible, low-rank attractor concentrated in a small set of attention heads. This both echoes the localization-head hypothesis---now induced by explicit supervision---and explains the method's data efficiency.

We are equally clear about what does \emph{not} happen. The large intrinsic gains do not yet translate into downstream benchmark gains: the aligned models track the base model, and saliency-map quality and task grounding are dissociated, a caution against treating sharper attention as a proxy for better answers. We also identified a clean failure mode---alignment withdraws attention from background ``stuff'' the annotations never cover---and its direct fix in stuff-aware supervision.

These findings set up the work in progress. Porting the objective into the standard LLaVA-1.5 instruction-tuning recipe with LoRA is designed to retain the intrinsic gains while preserving general capabilities and to test whether decoupling the signal from a caption-only diet closes the transfer gap. Several caveats bound the present claims: results are for a single architecture and a single $\lambda$, parameter drift is a proxy for functional change, and reproducibility is shown across training configurations but not yet across random seeds. Closing the transfer gap, and validating the candidate localization heads causally, are the natural next steps. We hope that treating attention as a directly supervisable signal, rather than only a diagnostic one, proves a useful lever for building VLMs that look where they should.
